\section{Conclusion}
\label{sec:conclusion}
This work contributes a reproducible reference implementation
for coupled Brinkman--convection-diffusion topology optimisation
in FEniCSx, including adjoint sensitivity verification,
automated testing, containerised execution,
and openly archived releases.

The framework combines coupled Brinkman--transport physics
and SUPG-stabilised continuous-adjoint optimisation
within a reusable FEniCSx implementation.
It further provides software validation, reproducibility infrastructure,
default benchmark configurations, automated testing,
and openly archived releases
(Tables~\ref{tab:adjoint_verif},~\ref{tab:repro_checklist};
DOI:~\href{https://doi.org/10.5281/zenodo.20538833}{10.5281/zenodo.20538833}).

On an $80\times20$ mesh, the framework achieves
RMSE~$=\rmseTopOpt$ for the linear gradient benchmark,
demonstrating successful execution of the integrated workflow
(Section~\ref{sec:taylor_test});
thresholding results are in Supplementary~S5.
The target profile gallery (Table~\ref{tab:gallery}) demonstrates
that the framework operates across multiple outlet objectives
by changing a single argument.

\paragraph{Future directions.}
The framework is designed as a reusable foundation for future research
on coupled flow-transport topology optimisation.
Principal open directions are:
(i)~robust projection~\cite{wang2011} integrated into the
optimisation loop;
(ii)~unstructured mesh support and larger-scale computations;
(iii)~higher Pe regimes where advective manipulation
provides greater benefit over passive diffusion.

Beyond the specific benchmark considered here,
\texttt{brinkgrad} is intended as a reusable reference implementation
for researchers developing related PDE-constrained optimisation
workflows in FEniCSx.

\texttt{brinkgrad} provides a practical, documented, and
reproducible reference implementation for coupled flow-transport
adjoint-based topology optimisation in FEniCSx.
By providing a fully tested, containerised, and documented baseline,
\texttt{brinkgrad} lowers the barrier to entry for researchers
developing advanced microfluidic and porous-media devices,
shifting the community focus from re-implementing baseline solvers
to innovating on physical models and objectives.
